\clearpage
\VerdierPhysicalPageAnchor{0208}%
\begingroup\small\setlength{\parskip}{0pt}\setlength{\emergencystretch}{1.5em}
\setlength{\abovedisplayskip}{4pt plus 2pt minus 2pt}%
\setlength{\belowdisplayskip}{4pt plus 2pt minus 2pt}%

\VerdierSemanticAnchor{ch3d-3}
\section*{3. \texorpdfstring{$\operatorname{Ext}$}{Ext} の研究。}

\VerdierSemanticAnchor{ch3d-3-preamble}
\noindent\textit{この節では、$\mathcal A$ はアーベル圏を表す。}

\VerdierSemanticAnchor{ch3d-3-1}
\subsection*{3.1. 種々の同値な定義。}

\VerdierSemanticAnchor{ch3d-3-1-1}
\noindent\textbf{3.1.1.}\quad
$\mathcal A$ の対象の任意の対 $X$ と $Y$ および任意の整数 $p>n$ に対して
$\operatorname{Ext}_{\mathcal A}^{p}(X,Y)=0$ となるとき（(1.2.5), (1.2.9)）、
\emph{$\mathcal A$ のコホモロジー次元は $n$ 以下である}という。
$\mathcal A$ のコホモロジー次元が $n$ 以下となる最小の整数 $n$ を
\emph{$\mathcal A$ のコホモロジー次元}と呼び、
$\operatorname{dimcoh}(\mathcal A)$ と記す。$\mathcal A$ の
コホモロジー次元が $n$ 以下となるような整数 $n$ が存在しないとき、
$\mathcal A$ のコホモロジー次元は無限であるといい、
$\operatorname{dimcoh}(\mathcal A)=+\infty$ と書く。

\VerdierSemanticAnchor{ch3d-3-1-2}
\noindent\textbf{3.1.2.}\quad
(1.1.3) を思い出すと、$\operatorname{Qis}^{*}$ は
$K^{*}(\mathcal A)$ の擬同型全体を表す
（$*=b,-,+,{}$「空」）。$X$ を $K^{*}(\mathcal A)$ の対象とする。
$\operatorname{Qis}^{*}/X$（それぞれ
$X\backslash\operatorname{Qis}^{*}$）を、終域が $X$（それぞれ始域が
$X$）である擬同型の圏とする（第 II 章，2.2.1）。
$\operatorname{Qis}^{*}$ は射の乗法的系なので、圏
$\operatorname{Qis}^{*}/X$（それぞれ
$X\backslash\operatorname{Qis}^{*}$）は左フィルター圏（それぞれ
右フィルター圏）である（第 II 章，2.2.1）。圏
$\operatorname{Qis}^{*}/X$（それぞれ
$X\backslash\operatorname{Qis}^{*}$）が始対象（それぞれ終対象）
$u\colon Y\longrightarrow X$（それぞれ $u\colon X\longrightarrow Y$）
をもつための必要十分条件は、$Q^{*}$-左自由（それぞれ右自由）な対象
$Y$ と擬同型 $u\colon Y\longrightarrow X$（それぞれ
$u\colon X\longrightarrow Y$）が存在することである。これは
（第 II 章，2.3.3）を用いれば直ちに分かる。

\VerdierSemanticAnchor{ch3d-3-1-3}
\medskip\noindent\textbf{命題 3.1.3.}\quad
\textit{$X$ と $Y$ を $\mathcal A$ の二つの複体とする。}

\VerdierSemanticAnchor{ch3d-3-1-3-1}
\[
\textit{(1)}\qquad
\operatorname{Ext}_{\mathcal A}^{0}(X,Y)=
\varinjlim_{(\operatorname{Qis}/X)^o}
\operatorname{Hom}_{K(\mathcal A)}(\mathord\cdot,Y).
\]

\VerdierSemanticAnchor{ch3d-3-1-3-2}
\[
\textit{(2)}\qquad
\operatorname{Ext}_{\mathcal A}^{0}(X,Y)\simeq
\varinjlim_{Y\backslash\operatorname{Qis}}
\operatorname{Hom}_{K(\mathcal A)}(X,\mathord\cdot).
\]

\VerdierSemanticAnchor{ch3d-3-1-3-3}
\[
\textit{(3)}\qquad
\operatorname{Ext}_{\mathcal A}^{0}(X,Y)\simeq
\varinjlim_{(\operatorname{Qis}/X)^o\times
  (Y\backslash\operatorname{Qis})}
\operatorname{Hom}_{K(\mathcal A)}(\mathord\cdot,\mathord\cdot).
\]

\VerdierSemanticAnchor{ch3d-3-1-3-4}
\[
\textit{(4)}\qquad
\operatorname{Ext}_{\mathcal A}^{n}(X,Y)\simeq
\operatorname{Ext}_{\mathcal A}^{0}(X[p],Y[n+p]),
\qquad\textit{任意の整数 $p$ に対して。}
\]

\VerdierSemanticAnchor{ch3d-3-1-3-5}
\[
\textit{(5) $K^{-}(\mathcal A)$ の任意の $X$ に対して、}\qquad
\operatorname{Ext}_{\mathcal A}^{0}(X,Y)\simeq
\varinjlim_{(\operatorname{Qis}^{-}/X)^o}
\operatorname{Hom}_{K(\mathcal A)}(\mathord\cdot,Y).
\]

\VerdierSemanticAnchor{ch3d-3-1-3-6}
\[
\textit{(6) $K^{+}(\mathcal A)$ の任意の $Y$ に対して、}\qquad
\operatorname{Ext}_{\mathcal A}^{0}(X,Y)\simeq
\varinjlim_{Y\backslash\operatorname{Qis}^{+}}
\operatorname{Hom}_{K(\mathcal A)}(X,\mathord\cdot).
\]
\endgroup
